\section{Objetivo e Questões de Pesquisa}

\textbf{GQM.} Analisar o uso de IA generativa aplicada à programação (GitHub Copilot, ChatGPT, Claude, Cursor etc.) com o propósito de caracterizar e avaliar com respeito a velocidade percebida nas entregas e qualidade percebida do código produzido, sob o ponto de vista de desenvolvedores de software profissionais, no contexto de profissionais atuantes no Brasil (empresas, consultorias, setor público e autônomos).

\begin{itemize}
  \item \textbf{QP1:} Em que medida desenvolvedores percebem ganhos de velocidade e produtividade ao usar assistentes de IA? (descritiva)
  \item \textbf{QP2:} Em que medida percebem impactos positivos ou negativos na manutenibilidade, legibilidade e taxa de defeitos do código assistido por IA? (descritiva)
  \item \textbf{QP3:} Há diferença estatisticamente significativa na percepção de velocidade e de qualidade entre níveis de senioridade (Júnior, Pleno, Sênior)? $H_0$: as distribuições dos escores são iguais nos três níveis; $H_1$: ao menos um nível difere. (confirmatória)
  \item \textbf{QP4:} Há correlação entre o rigor com que o desenvolvedor inspeciona/testa o código gerado por IA e a qualidade percebida do código assistido por IA? $H_0$: $\rho = 0$; $H_1$: $\rho \neq 0$. (confirmatória)
\end{itemize}

Toda análise fora das QP3 e QP4 (e.g., por ferramenta ou frequência de uso) é exploratória e será rotulada assim.
